\section{Formal setting}
\label{sec:formal-setting}

This section fixes the formal framework used throughout the paper.
All symbols, objects, and relations are taken exclusively from the
Ontology of Continua (OC Core) and the executable specification
ETS.K\_EA.v1.1. No extensions, reinterpretations, auxiliary formalisms,
or external semantics are introduced.

\subsection{Enterprise continuum}
\label{subsec:enterprise-continuum}

An enterprise is modeled as a continuum $K_{EA}$ with the following
components, as defined in ETS.K\_EA.v1.1:
\[
K_{EA} =
\bigl(
\Omega_{EA},
\partial\Omega_{EA},
A_{EA},
P_{EA},
\Theta_{EA},
J_{EA},
C_{EA},
k_{EA},
\tau_{EA}
\bigr).
\]

Here:
\begin{itemize}
\item $\Omega_{EA}$ is the internal state space of the enterprise
      continuum;
\item $\partial\Omega_{EA}$ denotes its boundary, including regions
      relevant for collapse and STOP conditions;
\item $A_{EA}$ is the set of admissible structural axes;
\item $P_{EA}$ denotes potentials associated with these axes;
\item $\Theta_{EA}$ is the set of admissible thresholds;
\item $J_{EA}$ represents internal flows;
\item $C_{EA}$ denotes admissible cycles;
\item $k_{EA}$ is the continuity measure;
\item $\tau_{EA}$ is the internal time parameter.
\end{itemize}

No assumptions are made beyond those explicitly fixed in the executable
specification.

\subsection{Admissible phase logic}
\label{subsec:phase-logic}

The enterprise continuum admits the following phase outcomes, treated
as mutually exclusive state predicates over $\Omega_{EA}$:
\[
\text{OK}, \quad \text{INERTIA}, \quad \text{COLLAPSE}, \quad \text{STOP}.
\]

These phases are defined in ETS in terms of threshold relations,
boundary interactions, and structural conditions. In particular,
$\text{STOP}$ denotes a terminal or boundary condition at which further
meaningful state discrimination within the enterprise continuum is no
longer admissible.

Throughout this paper, phase outcomes are treated purely as logical
predicates. No dynamic transition rules, recovery mechanisms, temporal
ordering of phases, or intervention semantics are assumed or implied.

\subsection{Observability constructs}
\label{subsec:observability-constructs}

Observability is restricted to the constructs already present in ETS.
Let:
\begin{itemize}
\item $\Omega_{\mathrm{obs}} \subseteq \Omega_{EA}$ denote the observable
      projection of the enterprise state space;
\item $\Sigma_{\mathrm{obs}}$ denote the space of admissible observable
      signatures;
\item $\Pi$ denote the observation operator
      $\Pi : \Omega_{EA} \rightarrow \Sigma_{\mathrm{obs}}$.
\end{itemize}

The operator $\Pi$ captures all admissible observations. No additional
measurement channels, inferred quantities, reconstructed states, or
derived observables are permitted. In particular, $\Pi$ is not assumed
to be injective: distinct internal states may induce identical
observable signatures.

\subsection{Observational consistency}
\label{subsec:observational-consistency}

For a given observable signature
$\sigma \in \Sigma_{\mathrm{obs}}$, the set of internally compatible
enterprise states is given by:
\[
\Pi^{-1}(\sigma) \subseteq \Omega_{EA}.
\]

This set may contain multiple, structurally distinct internal states or
even entire enterprise continua that are observationally
indistinguishable. All results in this paper are formulated in terms of
invariance or non-invariance of predicates over such sets.

\subsection{Global meta-constraints}
\label{subsec:meta-constraints}

The following constraints apply globally throughout the paper:
\begin{itemize}
\item No inference, estimation, or reconstruction procedures are used or
      implied.
\item No probabilistic or statistical assumptions are introduced.
\item No algorithmic or computational notions of decidability are
      employed; decidability is treated strictly as a logical property.
\item No prescriptions, interventions, control strategies, or
      governance interpretations are derived.
\end{itemize}

Within these constraints, the question of decidability is addressed
solely as a structural consequence of the OC/ETS formalism under
incomplete observability.
